\documentclass[11pt]{article}
\usepackage[margin=1in]{geometry}
\usepackage{amsmath,amssymb,amsthm}
\usepackage[hidelinks]{hyperref}
\newtheorem{theorem}{Theorem}
\newtheorem{lemma}{Lemma}
\newcommand{\R}{\mathbb R}
\newcommand{\F}{\mathcal F}
\newcommand{\norm}[1]{\lVert#1\rVert}
\title{A theta-cardinal bound for spectral convergence\newline in Suzuki's localized Weil spaces}
\author{Zetral Research Project}
\date{25 November 2026 --- candidate proof, version 0.2}
\begin{document}
\maketitle
\begin{abstract}
Assuming the Riemann hypothesis and the sampling and inverse-compression theorems of Suzuki, we derive a fixed-index upper bound of order $X^{5/2}e^{-2\pi X}$ for the difference between the exact localized positive differentiation eigenvalue and its limiting distinct zeta-zero ordinate, where $X=e^{2a}$. The argument combines divided-out theta kernels, compact truncation and a quadratic Rayleigh--Ritz estimate. No simplicity assumption is needed. Explicit constants are supplied. Five supporting scalar lemmas have been checked in Lean; the analytic and spectral proof has not been formally verified or independently reviewed. Originality is not established.
\end{abstract}
\section{Setting and statement}
Write $\xi_0(s)=s(s-1)\pi^{-s/2}\Gamma(s/2)\zeta(s)/2$ and $A(z)=\xi_0(1/2-iz)$. Suzuki uses $2\xi_0$; the zeros and the normalized cardinal functions are unchanged. Assume RH throughout. Let $\Gamma$ denote the set of distinct real ordinates, $m_\rho$ their multiplicities, and $0<\gamma_1<\gamma_2<\cdots$ the positive ordinates. In particular,
\[
M=\sum_{\rho\in\Gamma}m_\rho/\rho^2<\infty.
\]
Use $\F u(z)=\int_{\R}u(x)e^{izx}\,dx$. In the weighted sampling model of Suzuki \cite[Eq.~(2.8)]{suzuki}, the local form completion embeds isometrically as a closed subspace $M_a$ of $\ell^2(\Gamma,m)$, containing the sample vectors of $C_c^\infty(-a,a)$. Let $P_a$ project onto $M_a$. The global inverse operator $C$ is diagonal with entries $1/\rho$. Suzuki's Theorem 3.3 identifies the inverse of the phase-$\pi$ self-adjoint local differentiation extension with the compression of $C$ to $M_a$. Thus its positive eigenvalues $t_{a,n}$ are reciprocals of the decreasing positive eigenvalues $\mu_{a,n}$ of that compression. We import these results, rather than re-proving them.
\begin{theorem}
For each fixed $n$, there exist $K_n,a_n>0$ such that, for $a\ge a_n$ and $X=e^{2a}$,
\[
0\le t_{a,n}-\gamma_n\le K_n X^{5/2}e^{-2\pi X}.
\]
\end{theorem}
The constants may depend on the target zeros and their xi derivatives. The bound is not a uniform estimate as $n$ grows, nor a computable certificate for finite-dimensional numerical roots.
\section{A quadratic Ritz estimate}
\begin{lemma}
Let $C$ be compact self-adjoint, $L=\inf\sigma(C)$, and $E$ the span of its first $n$ positive eigenvectors, with corresponding eigenvalues $\lambda_1\ge\cdots\ge\lambda_n>0$. For an orthogonal projection $P$ onto a closed space $V$, put $\delta=\norm{(I-P)|_E}$. If $\delta<1$ and $(\lambda_n-L)\delta^2<\lambda_n$, the $n$th positive Ritz value $\mu_n$ exists and
\[
0\le\lambda_n-\mu_n\le(\lambda_n-L)\delta^2.
\]
\end{lemma}
\begin{proof}
Let $Q$ project onto $E$ and let $B=QP|_E$. Then $(1-\delta^2)I\le B\le I$ and $PE$ has dimension $n$. For $u\in E$ and $v=Pu$, one has $\norm{v}^2=\langle Bu,u\rangle$ and $\norm{(I-Q)v}^2=\langle(B-B^2)u,u\rangle\le\delta^2\norm{v}^2$. Since $E$ reduces $C$, the Rayleigh quotient on $PE$ is at least $\lambda_n-(\lambda_n-L)\delta^2>0$. Apply max--min on $PE\subset V$ for the lower bound and on the full space for the upper bound.
\end{proof}
This is a special case of established squared-angle Ritz estimates \cite{ritz}.
\section{Cardinal functions with theta decay}
The classical even smooth theta kernel \cite[Eqs.~(1.1)--(1.2), Lemma~3.1]{theta} satisfies $\F\Phi=A$ and, for $x>0$,
\[
\Phi(x)=\sum_{k\ge1}(4\pi^2k^4e^{9x/2}-6\pi k^2e^{5x/2})e^{-\pi k^2e^{2x}}.
\]
With $y=e^{2|x|}$, its series gives $|\Phi(x)|\ll y^{9/4}e^{-\pi y}$ by factoring out $e^{-\pi y}$ and bounding the remaining convergent series at $y=1$. For any fixed real $\alpha$ and $Y\ge1$,
\[
\int_Y^\infty v^{\alpha-1}e^{-\pi v}\,dv\ll_\alpha Y^{\alpha-1}e^{-\pi Y};
\]
substitute $v=Y+w$ and bound $(1+w/Y)^{\alpha-1}$ by $(1+w)^{\max(\alpha-1,0)}$.
Fix a real zero $\gamma$ of multiplicity $m$. Set $f_0=\Phi$ and, for $0\le r<m$, set
\[
f_{r+1}(x)=-i e^{-i\gamma x}\int_{-\infty}^x e^{i\gamma t}f_r(t)\,dt.
\]
Inductively, $\F f_r(z)=A(z)/(z-\gamma)^r$ and $|f_r(x)|\ll_{\gamma,r}y^{9/4-r}e^{-\pi y}$. Indeed, for $r<m$ the transform vanishes at $\gamma$, so the integral may be replaced by the negative outward right tail when $x\ge0$. The left tail is used when $x\le0$. The displayed tail estimate reduces the power of $y$ by one. Differentiating gives $(i\partial_x-\gamma)f_{r+1}=f_r$. Integration by parts with vanishing endpoints proves the transform identity, first on the real axis and then by entire continuation; the decay ensures entire transforms. This closes the induction.

Normalize $q_\gamma=f_m/(\sqrt m\,A^{(m)}(\gamma)/m!)$. Its sample vector is the unit coordinate vector $e_\gamma=m^{-1/2}\mathbf1_{\{\gamma\}}$. The quotient has removed exactly the multiplicity at $\gamma$ and retains all other zeros. We have
\[
|q_\gamma(x)|\ll_\gamma y^{9/4-m}e^{-\pi y},\qquad
|q'_\gamma(x)|\ll_\gamma y^{13/4-m}e^{-\pi y}.
\]
The derivative estimate follows from the differential identity above. No local space membership of the untruncated kernel is needed.
\section{Truncation and proof of the theorem}
Choose a smooth transition $\eta$ equal to zero on $(-\infty,1/2]$ and one on $[1,\infty)$, with values in $[0,1]$. For $a\ge1$, put $h=X^{-1}$ and $\chi_a(x)=\eta((a-|x|)/h)$. Then $u=\chi_a q_\gamma\in C_c^\infty(-a,a)$ and the omitted tail is $r=(1-\chi_a)q_\gamma$. The inner edge of the transition satisfies $Y=e^{2(a-h)}=Xe^{-2/X}\ge X-2$. Hence powers of $Y$ and $X$ are comparable, and $e^{-\pi Y}\le e^{2\pi}e^{-\pi X}$.

The product rule gives
\[
\norm{r'}_1\le\int_{|x|\ge a-h}|q'_\gamma(x)|\,dx+\int|\chi'_a q_\gamma|\,dx
\ll_\gamma X^{9/4-m}e^{-\pi X}.
\]
The first integral uses the tail estimate. The second uses $|\chi'_a|\ll h^{-1}$ on strips of total length $h$ and the bound on $q_\gamma$. Integrating by parts at every nonzero ordinate gives $|\F r(\rho)|\le\norm{r'}_1/|\rho|$. Consequently, for the sample vector $v$ of $u$,
\[
\norm{v-e_\gamma}_{\ell^2(\Gamma,m)}^2\le M\norm{r'}_1^2
\ll_\gamma X^{9/2-2m}e^{-2\pi X}\le C_\gamma X^{5/2}e^{-2\pi X}.
\]
Projection minimizes distance. For $E=\operatorname{span}(e_{\gamma_1},\ldots,e_{\gamma_n})$, Cauchy--Schwarz therefore gives $\delta_a^2=\norm{(I-P_a)|_E}^2\ll_n X^{5/2}e^{-2\pi X}$. In the lemma take $\lambda_n=1/\gamma_n$ and $L=-1/\gamma_1$. Put $B_n=\lambda_n-L$. For sufficiently large $a$, $B_n\delta_a^2\le\lambda_n/2$. Reciprocating the Ritz bounds yields
\[
0\le t_{a,n}-\gamma_n\le 2B_n\gamma_n^2\delta_a^2,
\]
which proves the theorem under the stated imports and RH.
\section{Explicit constants}
The theta series gives $|\Phi(x)|\le128y^{9/4}e^{-\pi y}$. Indeed $k\le2^{k-1}$ and $k^2-1\ge3(k-1)$ for integers $k\ge1$, so $k^4e^{-\pi(k^2-1)}\le4^{-(k-1)}$. Summing and using $4\pi^2+6\pi<88$ proves the stated bound. For $\alpha\le9/4$ and $Y\ge1$,
\[
\int_Y^\infty v^{\alpha-1}e^{-\pi v}\,dv
\le\frac{Y^{\alpha-1}e^{-\pi Y}}{\pi-5/4}.
\]
This follows from $v=Y+w$ and $(1+w/Y)^{\alpha-1}\le e^{5w/4}$. Each outward integration in the divided-kernel recurrence contributes a factor at most $1/[2(\pi-5/4)]<1/3$. Write
\[
d_\gamma=|A^{(m)}(\gamma)|/m!,\qquad D_\gamma=\frac{128}{3^m\sqrt m\,d_\gamma}.
\]
Then $|q_\gamma|\le D_\gamma y^{5/4}e^{-\pi y}$ and $|q'_\gamma|\le D_\gamma(|\gamma|+3)y^{9/4}e^{-\pi y}$, since $m\ge1$. Choose the transition $\eta$ nondecreasing, so $\int|\chi'_a|=2$. The function $y^{5/4}e^{-\pi y}$ decreases for $y\ge1$. Applying the tail bound at $Y\ge X-2$ to the derivative and using the total variation bound on the transition gives
\[
\norm{r'}_1\le H_\gamma X^{5/4}e^{-\pi X},\qquad
H_\gamma=D_\gamma e^{2\pi}\left(2+\frac{|\gamma|+3}{\pi-5/4}\right).
\]
The paired Hadamard product of the even order-one function $A$ gives $A(z)/A(0)=\prod_{\gamma>0}(1-z^2/\gamma^2)^{m_\gamma}$. Its convergence on compact sets follows from reciprocal-square summability; evenness removes the linear exponential. Differentiation at zero gives $M=-A''(0)/A(0)$. Set
\[
S_n=M\sum_{j=1}^n H_{\gamma_j}^2,\qquad K_n=2B_n\gamma_n^2 S_n.
\]
The preceding projection argument gives the theorem with this $K_n$ whenever
\[
X\ge e^2,\qquad B_n S_n X^{5/2}e^{-2\pi X}\le1/(2\gamma_n).
\] One sufficient threshold is $X\ge\max(e^2,\log(2\gamma_n B_n S_n)/\pi)$, using $X^{5/2}\le e^{\pi X}$ for $X\ge1$. The constants require the ordinates, exact multiplicities and nonzero derivatives, but not an infinite zero list. Certified enclosures for these inputs would make the estimate effective for a specified index; finite-prime-only computability is not established.
\section{Scope, computation, and review status}
The bound concerns the exact phase-$\pi$ differentiation extension, not the ordinary $L^2$ Weil-form minimum or a finite Fourier root. It asserts neither optimality nor uniformity in the index. The reference calculations reproduce the reported rounded tables; Rust 1.85.0 passes three baseline tests and the six-root smoke test at $X=3,N=18$. At $X=11,n=1$, a non-certified 50--60 digit evaluation of the explicit constant estimates the upper bound as $3.372\times10^{-11}$. The ball-arithmetic script is staged in draft PR 2, but its execution was not verified. Neither this estimate nor numerical agreement is a certified enclosure, and finite Fourier approximation, quadrature and root-isolation errors remain uncontrolled by the theorem.

The divided entire-function basis in Suzuki's earlier Weil-space work \cite{suzuki2025}, Fourier interpolation \cite{brs}, and squared-angle Ritz estimates \cite{ritz} are antecedents. Rapid-window estimates for the bottom of the Weil form, including Zhu \cite{zhu} and the reported bounds of Mori \cite{mori}, are close in technique but concern different spectral quantities. The spectral determinant construction of Connes--Consani--Moscovici \cite{ccm} and the numerical realization of Kim et al. \cite{kim} also differ from the specific exact phase-$\pi$ cutoff rate. The repository's targeted prior-art audit did not identify an equivalent theorem in the inspected material; incomplete full-text access and the absence of exhaustive citation-index review preclude any priority claim.

Five scalar inequalities supporting the variational and reciprocal steps were kernel-checked with Lean 4.19.0 and pinned Mathlib in two CI runs. Their axiom reports contain only \texttt{propext}, \texttt{Classical.choice}, and \texttt{Quot.sound}. The geometric, analytic and operator-theoretic premises remain unformalized. This note was generated and internally reviewed with AI assistance under human research direction; independent mathematical review remains necessary before academic submission. The detailed proofs, audits, logs and data are at \url{https://github.com/sguzman/zetral}.
\begin{thebibliography}{9}
\small
\setlength{\itemsep}{0pt}
\bibitem{suzuki} M. Suzuki, \emph{Spectral convergence associated with Weil's quadratic form}, arXiv:2610.32704v1 (2026).
\bibitem{theta} J. C. Lagarias and D. Montague, \emph{The integral of the Riemann xi-function}, arXiv:1106.4348v3 (2011).
\bibitem{ritz} A. V. Knyazev and M. E. Argentati, \emph{Rayleigh--Ritz majorization error bounds with applications to FEM}, SIAM J. Matrix Anal. Appl. 31 (2010), doi:10.1137/08072574X.
\bibitem{suzuki2025} M. Suzuki, \emph{On the Hilbert space derived from the Weil distribution}, Canadian Journal of Mathematics, published online 2025, doi:10.4153/S0008414X25101739.
\bibitem{brs} A. Bondarenko, D. Radchenko, K. Seip, \emph{Fourier interpolation with zeros of zeta and L-functions}, Constructive Approximation 57 (2023), doi:10.1007/s00365-022-09599-w.
\bibitem{zhu} X. Zhu, \emph{Weil positivity in compact windows}, arXiv:2608.24827v2 (2026).
\bibitem{mori} R. Mori, \emph{Unconditional doubly exponential upper bounds for windowed Weil quadratic forms}, repository summary (accessed 25 November 2026), \url{https://github.com/moriryota/weil-window-upper-bounds}.
\bibitem{ccm} A. Connes, C. Consani, H. Moscovici, \emph{Zeta spectral triples}, arXiv:2511.22755v1 (2025).
\bibitem{kim} T. Kim et al., \emph{A numerical realization of Suzuki's Weil-quadratic-form operator}, arXiv:2607.24830v3 (2026).
\end{thebibliography}
\end{document}
